

\input{settings_FM_template.tex}


\usepackage{palatino}
\usepackage{graphicx}
\usepackage{transparent}

\setbeamercovered{transparent=35}

\title[Editor konfiguračních souborů~Flow123d]{Editor konfiguračních souborů~Flow123d}
\subtitle{Diplomová práce}
\author[Bc. Tomáš Křížek]{Bc. Tomáš Křížek}
\institute[TUL]{Technická univerzita v Liberci}
\date{19.~dubna~2016}
\newcommand{\TextTitulniStranaPodLinkou}{\tiny
Studentská 2 {\color{FM_TUL} |} 461\,17 Liberec 2 {\color{FM_TUL} |} {tomas.krizek1@tul.cz} {\color{FM_TUL} |} 
\href{http://www.fm.tul.cz/}{www.fm.tul.cz}}

\begin{document}
%\setbeamertemplate{caption}{\insertcaption}


\begin{frame}
	\titlepage
\end{frame}

\begin{frame}
	\frametitle{Obsah}
	\begin{itemize}
		\item Problematika
		\item Motivace
		\item Zpracování konfiguračního souboru
		\item Validace
		\item Kontextová nápověda
		\item Grafické rozhraní
		\item Použité technologie
	\end{itemize}
\end{frame}


\begin{frame}[t]
	\frametitle{Problematika}
	\begin{block}{Simulátor Flow123d}
	\begin{itemize}[<+->]
		\item modelování procesů v horninovém prostředí
		\item pouze textové rozhraní
		\item výpočetně náročné $\rightarrow$ vzdálené spouštění
	\end{itemize}
	\end{block}
	
	\begin{block}{Konfigurační soubory}
	\begin{itemize}[<+->]
		\item definice úlohy
		\item rozsáhlé možnosti nastavení
		\item záznamy, pole, primitivní datové typy, abstraktní záznamy
	\end{itemize}
	\end{block}
\end{frame}

\begin{frame}[t]
	\frametitle{Motivace}
	\begin{block}{Chyby v konfiguračních souborech}
	\begin{itemize}[<+->]
		\item časově náročná detekce a odstranění
		\item uživatelsky nepřívětivé
	\end{itemize}
	\end{block}
	
	\begin{block}{Vytváření a upravování konfiguračních souborů}
	\begin{itemize}[<+->]
		\item běžné textové editory
		\item rozsáhlá referenční dokumentace formátu
	\end{itemize}
	\end{block}
	
	\onslide<5>{
	\vspace{0.5cm}
	$\Rightarrow$ \textit{specializovaný editor}
	}
\end{frame}

\begin{frame}[fragile]
	\frametitle{Zpracování konfiguračního souboru}
	\only<1>{
	\hspace*{-5pt}\includegraphics[width=\textwidth]{../../img/data_structure_chain_1.pdf}\\
	}
	\only<2-5>{
	\hspace*{-5pt}\includegraphics[width=\textwidth]{../../img/data_structure_chain_2.pdf}\\
	}
	\only<6-9>{
	\hspace*{-5pt}\includegraphics[width=\textwidth]{../../img/data_structure_chain_3.pdf}\\
	}
	\only<10>{
	\hspace*{-5pt}\includegraphics[width=\textwidth]{../../img/data_structure_chain_4.pdf}\\
	}
	\begin{minipage}[t]{0.45\textwidth}
	\begin{block}{Formát YAML}
	\begin{itemize}
		\item<3-> nahrazuje formát CON
		\item<4-> přehledná syntaxe
		\item<5-> návrh použití syntaxe pro zápis (abstraktní záznamy, reference)
		%\item<6-> vlastní implementace zpracování souboru
	\end{itemize}
	\end{block}
	\end{minipage}
	\hspace{5pt}
	\begin{minipage}[t]{0.45\textwidth}
	\begin{block}{Autokonverze}
	\begin{itemize}
		\item<7-> speciální zkrácený zápis polí či záznamů
		\item<8-> v praxi -- jiný datový typ, než se očekává
		\item<9-> libovolně vnořené $\rightarrow$ rekurzivní průchod datové struktury
	\end{itemize}
	\end{block}
	\end{minipage}
\end{frame}

\begin{frame}
	\frametitle{Validace}
	\begin{itemize}[<+->]
	\item včasné odhalení chyb
	\item přizpůsobí se verzi Flow123d
	\item upozornění na možné problémy
	\end{itemize}
	\only<1-3>{
		\includegraphics[width=\textwidth]{img/validation_pause.png}
	}
	\only<4->{
		\includegraphics[width=\textwidth]{img/validation.png}
	}
\end{frame}


\begin{frame}
	\frametitle{Kontextová nápověda}
	\begin{itemize}[<+->]
	\item alternativa k rozsáhlé referenční dokumentaci
	\item zobrazuje relevantní dokumentaci k právě editované hodnotě (pozici)
	\item možnost navigace v rámci nápovědy
	\end{itemize}
	\only<1-3>{
		\includegraphics[width=\textwidth]{img/gui_doc_mortar_method_pause.png}
	}
	\only<4->{
		\includegraphics[width=\textwidth]{img/gui_doc_mortar_method.png}
	}
\end{frame}

\begin{frame}
	\frametitle{Grafické rozhraní aplikace}
	\includegraphics[height=7cm]{img/gui.png}
\end{frame}

\begin{frame}
	\frametitle{Použité technologie}
	\begin{itemize}[<+->]
	\item Python 3.4
	\item grafická knihovna PyQt 5
	\item knihovna pro textový editor QScintilla 2
	\item knihovna pro YAML -- PyYAML
	\item instalační balíčky pro Windows, Debian
	\end{itemize}
\end{frame}

\begin{frame}
	\frametitle{Shrnutí -- funkce editoru}
	\begin{itemize}[<+->]
	\item zpracování YAML
	\item autokonverze
	\item validace
	\item kontextová nápověda
	\item automatické doplňování textu
	\item vizualizace datové struktury
	\item přizpůsobí se verzi Flow123d
	\item multiplatformní aplikace (Windows XP a novější, Linux)
	\end{itemize}
\end{frame}

\begin{frame}{}{}
\begin{center}
\huge Děkuji za pozornost.
\end{center}
\end{frame}


\end{document}
